\chapter{Background \& Related Work}\label{background-related-work}

\section{Musical Metacreation and Machine Musicianship}\label{musical-metacreation-and-machine-musicianship}

The term Musical Metacreation \citep{pasquier2016introduction} provides the broad
research context describing systems that use computation to participate
in musical tasks such as composition, performance, improvisation and
accompaniment. However, machine creativity within this field is not
uniformly defined. Some systems frame creativity as the autonomous
production of musical artefacts, while others understand it as an
interactive process shaped by human input, shared agency and real-time
response.

This section reviews the literature across three conceptual axes:
autonomous generation versus co-creative interaction, machine output
versus human-machine process and independent musical authorship versus
responsive machine musicianship. This framing is important for
conceptualising the essence of creativity within IT4s.

\subsection{Musical Metacreation As A Field}\label{musical-metacreation-as-a-field}

Musical metacreation, often abbreviated as MuMe, describes computational
systems that take on creative musical tasks. The definition's breadth is
useful for situating Intelli-Trading Fours but requires a more precise
classification of the project's contribution. Rather than treating all
AI music systems as a homogenous category, the following discussion
distinguishes between systems that generate musical material
independently and systems whose creative behaviour is contingent on live
human interaction.

For IT4s, the usefulness of MuMe is not that it legitimises any
computational music system as ``creative'', but that it provides a
vocabulary for locating different forms of machine musical agency. The
project sits inside this field because it delegates parts of musical
interaction to software: listening, representation, analysis, planning
and playback. However, its contribution is narrower than autonomous
musical authorship. The system's musical behaviour depends on a human
rhythmic turn and is therefore better understood as interactive machine
musicianship.

\subsection{Autonomous Generative Agents}\label{autonomous-generative-agents}

One strand of musical metacreation frames machine creativity primarily
through autonomous generation. In such systems, creative behaviour
emerges from processes that are largely self-contained within the
computational domain through learned models, internal rules or agent
based interaction with other software entities.

Eigenfeldt's Musebots exemplify this approach, demonstrating musical
collaboration between autonomous agents whose creative agency remains
largely within the computational system without strict temporal
dependency on external input rather than being negotiated in an exchange
with a live performer \citep{eigenfeldt2013musebots}. While this model is useful for
understanding machine autonomy, it differs from the interactional
problem addressed in this work, where responses must be derived from
temporally bounded human input.

\subsection{Interactive And Co-Creative Machine Musicianship}\label{interactive-and-co-creative-machine-musicianship}

A contrasting perspective frames creativity not solely within autonomous
generation, but within the interaction between a computational system
and a human performer. In this model, the machine's musical behaviour is
informed by what the performer has just played, positioning the system's
contribution as a response, continuation, transformation or
accompaniment rather than as an isolated artefact.

Pachet's Continuator is a significant example of this
performer-conditioned approach. The system learns from a musician's
input and generates stylistically related continuations in real time,
bridging autonomous generation with interactive musical performance
\citep{pachet2002continuator}. Its relevance lies less in its specific modelling
method, instead framing machine musicianship as contextually responsive
participation.

The Continuator demonstrates that a computational system can derive
musical material from a performer's contribution and return output that
is meaningful within an ongoing exchange, however, Intelli-Trading Fours
takes a different design route. Rather than pursuing style-learning
continuation, its architecture is explicitly segmented. This results in
a system more limited as a generative model, but more transparent and
interpretable as an interaction architecture.

\subsection{Product-Centred Versus Process-Centred Creativity}\label{product-centred-versus-process-centred-creativity}

The distinction between autonomous and interactive systems also affects
how creativity should be evaluated. Product-centred approaches focus on
the quality of the generated musical artefact, which is appropriate for
systems that produce self-contained compositions or standalone musical
outputs. However, interactive music systems require a more
process-centred view. Their creative value cannot be judged only by the
standalone quality of generated material when it also depends on how the
system participates in the unfolding exchange with a performer.

Smith's account of human-AI musical partnerships frames agency, control
and collaboration as interactional concerns, while Ostermann et al.'s
evaluation of reactive jazz accompaniment demonstrates that perceived
creativity and musical usefulness can be assessed through the user's
experience of real-time system behaviour rather than through output
alone \citep{smith2022humanai,ostermann2021evaluating}.

This is particularly relevant to Intelli-Trading Fours, where the
central contribution lies not in claiming a finished autonomous
generative model, but in an observable and inspectable closed-loop
pipeline.

\section{Human-AI Improvisation and Co-Creative Systems}\label{human-ai-improvisation-and-co-creative-systems}

Having positioned Intelli-Trading Fours as a form of interactive machine
musicianship, this section considers what makes human-AI musical
interaction meaningful in practice. In improvisational contexts, the
system's value does not depend only on the material it generates, but on
how it listens, responds, shares agency and sustains the performer's
sense of an unfolding musical exchange.

\subsection{Improvisation As Perception And Response}\label{improvisation-as-perception-and-response}

A central problem in human-AI improvisation is that the system must
treat performer input as musical context rather than as arbitrary data.
Thom's work on real-time interactive improvisation is useful here
because it frames improvisational companionship as a time-sensitive
problem of perception and response, where the system must interpret live
musical material quickly enough to produce an appropriate continuation
or accompaniment \citep{thom2000bob}. Similarly, Pachet's Continuator
demonstrates how performer input can become the basis for real-time
musical continuation, showing that machine-generated material can
function as part of a live exchange rather than as a detached
composition \citep{pachet2002continuator}. Together, these systems frame improvisation
as a perception-response loop in which listening and generation are
mutually dependent rather than separate stages.

Intelli-Trading Fours adopts this perception-response principle, but
narrows it to a turn-based rhythmic setting. This differs from systems
that attempt continuous accompaniment or stylistic continuation across
an open-ended performance. The specific musical basis for this
turn-taking model is developed further in Section 2.5.

\section{Reactive, Predictive and Generative Musical Agents}\label{reactive-predictive-and-generative-musical-agents}

The previous section established human-AI improvisation as an
interaction loop in which perception, timing and agency shape the
musical exchange. This raises a further question: what is the nature of
this interaction? Interactive music systems can be positioned along a
spectrum from reactive systems, which respond directly to current input;
to predictive systems, which model likely future behaviour; and
generative systems, which produce more independent or stylistically
developed musical material. This distinction is not hierarchical, but
reflects different approaches to modelling interaction, agency and
musical structure.

\subsection{Reactive Musical Interaction}\label{reactive-musical-interaction}

Reactive musical systems respond to performer input, environmental
events or changing musical state. Petit and Serrano's work on reactive
programming for interactive structured music is helpful here for
presenting how musical systems can combine deterministic structure with
live or nondeterministic interaction through explicit states, events and
patterns \citep{petit2021reactive}. This supports the view that
real-time musical responsiveness often depends on architecture as much
as on generation.

Systems adopting this approach often structure behaviour around
externally conditioned responses, where output is derived from recent or
current input rather than from an independent generative process. Such
systems may still incorporate structured decision-making and internal
constraints, but their defining behaviour is anchored in the interaction
context.

\subsection{Predictive Interactive Systems}\label{predictive-interactive-systems}

Predictive musical systems extend this idea by modelling likely future
behaviour. Martin and Torresen's Interactive Musical Prediction System,
for example, uses mixture density recurrent neural networks to predict
future user control data in real time \citep{martin2019interactive}. Such
systems can support powerful forms of musical assistance and
co-performance but they also introduce different trade-offs around
training data, opacity, evaluation and control.

\subsection{Generative Response And Drum-Performance Models}\label{generative-response-and-drum-performance-models}

Generative musical agents prioritise the production of novel musical
material, often using learned models or rule-based systems to construct
patterns that extend beyond direct transformation of input.

Pachet's Continuator remains relevant here as it combines interaction
with stylistically informed continuation \citep{pachet2002continuator}. More recent
learning-based work also frames musical response generation as a
data-driven problem. Hu et al., for example, approach call-and-response
as automatic composition using a knowledge-enhanced model, while Gillick
et al. demonstrate learned drum-performance generation and humanisation
through sequence transformation \citep{hu2024responding,gillick2019learning}. These systems are useful comparison points because they
demonstrate response generation beyond simple transformation or
rule-based planning.

\subsection{Positioning Intelli-Trading Fours}\label{positioning-intelli-trading-fours}

The distinction between reactive, predictive and generative agents is
therefore not a hierarchy in which reactive systems are automatically
weaker. Vaggione's critique of simple deterministic-versus-random
framings is useful here because it suggests that musical structure can
emerge through constraints operating across different temporal levels
rather than through unconstrained generation alone \citep{vaggione1993determinism}.

Intelli-Trading Fours is best understood as a reactive, planning-led,
turn-based musical agent with partial generative structure. Its
behaviour is conditioned on recent human input, mediated through
interpretable analysis and planning stages. Compared with predictive and
fully generative systems, this approach prioritises transparency,
controllability and architectural clarity over expressive autonomy.

This partial generative structure reflects the core contribution of the
work: not autonomous musical creativity, but the design of a structured
and interpretable interaction system.

\section{Rhythm, Timing, Entrainment \& Turn-Taking}\label{rhythm-timing-entrainment-turn-taking}

It is important to note that the interaction problem is specifically
rhythmic. Rhythm is not simply musical content that can be generated
after the fact; it is a temporal medium through which performers
coordinate action, perceive structure and experience responsiveness
\citep{longuetHiggins1984rhythmic}. This makes timing a central design
concern for IT4s as responses must be situated appropriately within a
shared temporal framework.

\subsection{Rhythm As A Time-Sensitive Interaction Medium}\label{rhythm-as-a-time-sensitive-interaction-medium}

In musical interaction, rhythmic information is not experienced solely
as an abstract pattern; it is felt through the timing of events with
respect to the performer's expectation of when future events should
occur. This is especially important in percussion-based interaction,
where pitch and harmony do not contribute to its evaluation as much as
musical properties such as: onset timing, density, accent and dynamic
energy.

This requires positioning rhythm as both data and interaction. The
system should not merely need to detect that a drum hit happened; it
needs to know where that hit occurred within a temporal frame. Timing is
therefore not a secondary implementation detail, but a central musical
concept that the system must prioritise.

\subsection{Machine Listening And Rhythmic Feature Extraction}\label{machine-listening-and-rhythmic-feature-extraction}

Interactive rhythmic systems also require a form of machine listening.
Scheirer's work on music-listening systems becomes relevant here because
its framing of computational listening as the extraction of musically
meaningful information such as tempo, segmentation and perceptual
orientation \citep{scheirer2000musiclistening}. Although IT4s does not attempt full
audio-scene understanding, it still faces a related problem: live
performance events must be transformed into features that a
computational agent can use for response planning.

\subsection{Quantisation, Expressive Timing and Symbolic Representation}\label{quantisation-expressive-timing-and-symbolic-representation}

A key challenge in rhythmic interaction is that performed timing is
rarely perfectly aligned to a grid. Longuet-Higgins and Lee argue that
rhythmic interpretation involves resolving ambiguity through assumptions
about regularity, meter and grouping \citep{longuetHiggins1984rhythmic}.
This is relevant to IT4s whose architecture converts temporally messy
live input into a representation that is stable enough for analysis and
response planning. Quantisation is therefore not just a technical
convenience; it is an interpretive act that decides how performed events
relate to an underlying rhythmic structure.

At the same time, quantisation introduces a trade-off. A rigid grid can
make analysis and generation tractable but can also flatten expressive
timing. The musical importance of expressive timing in drum patterns is
highlighted by Gillick et al.'s work on drum performance with regard to
generation and humanisation. In particular, the research emphasises the
relationship between a performed groove and an associated abstract
rhythmic structure \citep{gillick2019learning}. IT4s employs a clear and
interpretable symbolic representation for the later planning stages,
while retaining timing offset information within the representation of a
compiled turn. This positions the system as suitable for analysis of a
deterministic nature, while maintaining extensibility in the form of
richer expressive timing and humanisation for future work.

\subsection{Implications For Intelli-Trading Fours}\label{implications-for-intelli-trading-fours}

The literature on rhythm, machine listening and timing supports a
central design decision in IT4s, the system therefore perceives rhythm
not only as musical content, but as an architectural coordination
problem. Timing must be preserved, discretised, interpreted and
rescheduled across the interaction loop. While Intelli-Trading Fours
does not implement continuous entrainment, the next section builds on
these ideas by considering trading fours specifically, where turn-taking
is not only a technical convenience but a musical practice that gives
the interaction its dialogic form.

\section{Trading Fours and Turn-Based Musical Dialogue}\label{trading-fours-and-turn-based-musical-dialogue}

Trading fours serves as a useful musical model for IT4s because it
contains a clear alternation between listening and response. In jazz
practice, a performer does not simply produce an isolated phrase; they
contribute musical material that invites a subsequent reply. Hodson and
Donnay et al. use trading fours as a model of interactive improvisation
and frame it as a form of musical communication involving collaborative
and emergent exchange \citep{donnay2014neural,hodson2007interaction}. For IT4s,
this means that trading fours is not only a stylistic reference, it
provides a structured interaction model in which the meaning of a
response depends on its relation to the preceding turn.

This structure is not a simplification but an execution of real
improvisation, giving the system a practical form of musical
accountability. Each AI response is evaluated as an answer to a specific
human turn, rather than as an independent generated artefact. This
supports the dissertation\textquotesingle s wider distinction between
autonomous generation and responsive interaction: IT4s does not claim to
be a fully open-ended improviser, but it does implement the conditions
needed for repeated rhythmic dialogue. The turn boundary makes the
source material identifiable, the response process inspectable and the
re-entry point testable.

The limitation is that this design reduces the simultaneity of human
musical interplay. IT4s does not currently negotiate tempo with the
performer. Instead, it treats trading fours as a discrete
call-and-response structure in which musical dialogue is mediated by
state transitions. This trade-off is appropriate for the current
prototype because it makes the interaction observable and technically
defensible: responsiveness can be discussed in terms of capture
accuracy, response timing, phase correctness and loop closure, while
richer forms of overlapping improvisation remain viable future work.

\section{Metrical Hierarchy, Salience \& Rhythmic Structure}\label{metrical-hierarchy-salience-rhythmic-structure}

After a structured form of rhythmic input has been achieved, the
challenge becomes determining which elements of that structure are
musically significant. Rhythm, more than simply a sequence of discrete
events, is better described as a perceptual construct shaped by
hierarchical organisation, accent and expectation.

\subsection{Metrical Hierarchy and Rhythmic Interpretation}\label{metrical-hierarchy-and-rhythmic-interpretation}

Research in the human perception of rhythm demonstrates that listeners
do not interpret rhythmic sequences as flat patterns of events, but as
structured hierarchies organised around an underlying meter.
\citet{longuetHiggins1984rhythmic} argue that rhythmic interpretation
depends on assumptions about regularity, grouping and metrical
structure, such that identical event sequences may be perceived
differently depending on their inferred context.

Within this hierarchical framework, not all events contribute equally to
perceived structure. Certain events become points of emphasis or
structural anchors that draw attention based on their position,
intensity or contextual contrast. Salience may arise from multiple
factors, including dynamic accent, placement on metrically strong
positions, relative isolation or proximity to phrase boundaries.

From a computational perspective, this implies the need for models that
estimate the relative importance of events rather than treating all
onsets uniformly. Such models typically combine heuristic indicators of
salience, allowing systems to approximate which elements of a rhythmic
pattern are likely to carry structural weight.

This abstraction gives way to systematic reasoning about rhythm,
including the identification of recurring patterns, structural positions
and relationships between events. While such representations necessarily
simplify continuous performance, they provide a tractable framework for
both analysis and generation, particularly in systems operating under
real-time constraints.

\subsection{Implications for Interactive Musical Systems}\label{implications-for-interactive-musical-systems}

Taken together, these perspectives highlight that rhythmic interaction
requires systems to operate on structured representations that capture
both hierarchical organisation and perceptual salience. For interactive
systems, this is particularly important because responses must relate
meaningfully to the perceived structure of input rather than to its raw
event stream.

Accordingly, effective interaction requires models that can:

\begin{itemize}
\item
  identify structurally significant events
\item
  capture phrase-level patterns and activity contours
\item
  operate on representations that balance interpretability with
  computational tractability
\end{itemize}

These requirements directly motivate the use of feature-based rhythmic
analysis in interactive systems, where structural descriptors provide a
bridge between raw input and higher-level response behaviour.

\section{Evaluation of Interactive Music Systems}\label{evaluation-of-interactive-music-systems}

Given the interaction-focused framing established in earlier sections,
evaluation in Intelli-Trading Fours cannot be reduced to the quality of
generated output alone. Prior work suggests generative music systems do
not admit a single objective ground truth, instead they are evaluated
across dimensions such as coherence, variability and interactivity
\citep{briot2020deep}. This complicates the use of
conventional output-based metrics, particularly in systems where musical
material is produced in response to external input rather than in
isolation.

This challenge is amplified in interactive contexts due to the
aforementioned research by \citet{smith2022humanai} and \citet{ostermann2021evaluating} who argue
place creative value within interaction and evaluate reactive
accompaniment systems through metrics that depend on real-time behaviour
and contextual alignment.

As a result, evaluation in this project is framed in terms of system
behaviour within the interaction loop. Rather than treating generated
rhythms as standalone artefacts, the evaluation focuses on architectural
and behavioural properties aligned with the system's design objectives.
These include timing correctness, state consistency, responsiveness,
interpretability and structural response behaviour. This approach
ensures that evaluation reflects the system's contribution as a
closed-loop, real-time interaction architecture, rather than as a fully
autonomous generative model.
